\section{Implementation}
\label{sec:implementation}

The Metrics Recorder was implemented in Python in our CASCADE fork. It follows
the architecture from Fig.~\ref{fig:metrics-architecture} and uses only standard
library functionality such as \texttt{json}, \texttt{time}, and
\texttt{contextvars} for its core operation. It integrates with CASCADE's pinned
\texttt{openai}~2.31.0 client for LLM usage data and the existing
\texttt{docker}~7.1.0 execution layer. Table~\ref{tab:implementation-artifacts}
summarizes the main implementation artifacts.

\subsection{Recorder and Metrics Utility}
\label{sec:implementation-recorder}

The extension introduces an abstract recorder interface that separates
instrumentation from storage. Its \path{record_event()} operation receives a
named event and arbitrary fields; \path{write_summary()} finalizes the run. The
default file-backed implementation owns one run. On construction it creates a
run identifier, stores the pipeline configuration, initializes its aggregates,
and emits a \texttt{run\_started} event.

For every subsequent event, the recorder adds a UTC timestamp and run ID, merges
the inherited metric context, normalizes values for JSON serialization, and
updates its in-memory aggregates. It then appends the event as one JSON object to
\path{metrics.events.jsonl}. The aggregates cover event counts and failures,
elapsed time, LLM usage by model, component, and phase, and per-sample values.
At finalization, they are written to \path{metrics.summary.json}. Monetary cost
is deliberately left unset because pricing depends on provider and time.

The Metrics Utility module stores the active recorder and context in two
\texttt{ContextVar} instances. Its \texttt{metric\_context()}
helper adds sample, method, component, or phase information to nested calls;
\texttt{time\_block()} measures elapsed time with \texttt{time.perf\_counter()}
and records success or failure without suppressing exceptions. The utility also
normalizes both prompt and completion tokens as well as input and output tokens
used by OpenAI-compatible APIs. Calls become no-ops when no recorder is active.

\begin{table}[ht]
    \centering
    \caption{Main implementation artifacts relative to \texttt{src/cascade/}.}
    \label{tab:implementation-artifacts}
    \scriptsize
    \begin{tabular}{@{}p{0.43\textwidth}p{0.45\textwidth}@{}}
        \toprule
        \raggedright Artifact & \raggedright Purpose \tabularnewline
        \midrule
        \raggedright \path{metrics/BaseMetricsRecorder.py}
            & \raggedright Recorder interface \tabularnewline
        \raggedright \path{metrics/MetricsRecorder.py}
            & \raggedright JSONL persistence and summary aggregation \tabularnewline
        \raggedright \path{utils/Metrics.py}
            & \raggedright Context, timing, and usage normalization \tabularnewline
        \raggedright \path{PipelineFactory.py}, \path{Pipeline.py}
            & \raggedright Construction and run lifecycle \tabularnewline
        \raggedright \path{generation/Generation.py}, \path{generation/executor/OpenAICaller.py}
            & \raggedright Attribution and LLM-call metrics \tabularnewline
        \raggedright \path{analysis/executor/JavaExecutor.py}, \path{utils/DockerizedWrapper.py}
            & \raggedright Java and Docker execution timings \tabularnewline
        \bottomrule
    \end{tabular}
\end{table}

\subsection{Pipeline Integration and Instrumentation}
\label{sec:implementation-integration}

The CLI forwards the \texttt{-o} output directory to
\texttt{PipelineFactory.build()}. The factory either loads a recorder selected
by an optional \texttt{MetricsRecorder} configuration block or constructs the
default implementation, passing it the output path and configuration snapshot.
It then injects the recorder into the pipeline. At the start of
\texttt{Pipeline.execute()}, the recorder is installed as the active sink. A
\texttt{finally} block records \texttt{run\_finished}, writes the summary, and
restores the previous context on both successful and failed runs. Metrics write
errors are reported but do not terminate the analysis.

Instrumentation covers both LLM work and non-LLM overhead. The pipeline times
extraction, filtering, and analysis. The analysis implementations add setup,
per-method, and teardown measurements. \path{generation/Generation.py} wraps
generation and repair calls in sample and component context, while the test
generator refines individual generation phases. \path{OpenAICaller.py} records
one event for every attempt, including failed retries, with model, endpoint,
finish reason, elapsed time, and token usage. Finally, the Java executor and
Docker wrapper measure project preparation, builder/container operations, and
test execution.

\subsection{Configuration, Execution, and Verification}
\label{sec:implementation-operation}

The fork is installed into its virtual environment with
\texttt{./cascade.venv/bin/pip install .} Metrics recording is enabled by
default; a custom recorder can be selected by name, module path, and keyword
arguments in the pipeline configuration. CASCADE is then invoked as follows:

\begin{quote}
\small\texttt{CASCADE run -i <project> -o <output> -c <config.json>}
\end{quote}

LLM-backed runs require \texttt{OPENAI\_API\_KEY}, or corresponding credentials
and a base URL for an alternative OpenAI-compatible endpoint. The execution path
also requires a local Java runtime and a reachable Docker daemon; Maven commands
run inside the configured Docker image. The raw event stream and summary are
written directly to the selected output directory.

With this setup, every standard CASCADE invocation automatically produces both
detailed events and aggregate metrics without requiring changes to individual
pipeline components. These artifacts provide the measurement data used for the
model and resource comparisons in Sect.~\ref{sec:evaluation}.

\sectioncontributors{Alexander Hristov and Caspar Moritz Klein.}
